\documentclass[11pt]{article}

\usepackage{geometry}
\geometry{left=1.25in, right=1.25in}

\usepackage{graphicx}
\graphicspath{ {./media/} }

\usepackage{svg}
\svgpath{{./media/}}
\svgsetup{inkscapelatex=false}

\usepackage{subcaption}
\usepackage{cleveref}

\makeatletter
\renewcommand{\maketitle}{
    {\Large \bf \@title}
    \vspace{0.75em}

    {\normalsize \@author}
    \vspace{0.25em}

    {\normalsize \@date}
    \vspace{1em}
}
\makeatother

\begin{document}
\title{Self-Assembly Exercises}
\date{\today}
\author{Makise Kurisu}
\maketitle

\subsection*{Exercise 0}
There are 11 different glue label and glue strength combinations (5 for each glue label of
strength 1, 5 for each glue label of strength 2, and an additional 1 for the null glue). Each of the
four sides of a tile can hold any of the 11 different glue labels. Therefore, there are a total of
$11^4 = 14641$ unique tile types you can make.

\subsection*{Exercise 1}
The goal of this temperature-2 system is to produce a ternary counter. In this system, there are 12
tile types evenly split into four categories: \verb|Border|, \verb|Zero|, \verb|One|, and \verb|Two|
as shown in \cref{fig:exs-1}.

\begin{figure}[htbp]
    \centering
    \begin{subfigure}{.5\textwidth}
        \includesvg[width=.30\linewidth]{exercise_1-SEED}
        \includesvg[width=.30\linewidth]{exercise_1-BOTTOM}
        \includesvg[width=.30\linewidth]{exercise_1-RIGHT}
        \caption{\texttt{Border} tiles.}
        \label{fig:exs-1_border}
    \end{subfigure}%
    \begin{subfigure}{.5\textwidth}
        \includesvg[width=.30\linewidth]{exercise_1-00}
        \includesvg[width=.30\linewidth]{exercise_1-21}
        \includesvg[width=.30\linewidth]{exercise_1-12}
        \caption{\texttt{Zero} tiles.}
        \label{fig:exs-1_zero}
    \end{subfigure}
    \begin{subfigure}{.5\textwidth}
        \includesvg[width=.30\linewidth]{exercise_1-01}
        \includesvg[width=.30\linewidth]{exercise_1-10}
        \includesvg[width=.30\linewidth]{exercise_1-22}
        \caption{\texttt{One} tiles.}
        \label{fig:exs-1_one}
    \end{subfigure}%
    \begin{subfigure}{.5\textwidth}
        \includesvg[width=.30\linewidth]{exercise_1-02}
        \includesvg[width=.30\linewidth]{exercise_1-20}
        \includesvg[width=.30\linewidth]{exercise_1-11}
        \caption{\texttt{Two} tiles.}
        \label{fig:exs-1_two}
    \end{subfigure}
    \caption{Each category of tiles for exercise 1.}
    \label{fig:exs-1}
\end{figure}

The \verb|Border| category forms the bottom and right borders of the ternary counter where the
``trit'' tiles, tiles in the categories \verb|Zero|, \verb|One|, and \verb|Two| can grow off of.
The seed assembly contains a single \verb|SEED| tile with which the \verb|BOTTOM| and \verb|RIGHT|
tiles can grow off of. As shown in \cref{fig:exs-1_border}, the \verb|SEED| tile has a strength-2
glue of colors \verb|R| and \verb|B| on its north and west sides. This allows the \verb|BOTTOM| and
\verb|RIGHT| tiles to grow horizontally and vertically forever, respectively, since the two tiles
each also have a strength-2 glue going in their respective direction.

The trit tiles represent each value in base-3 according to their name e.g., each tile in the
\verb|Zero| tiles represent ternary 0. Each row of the counter, made up of trit tiles, represents a
value in base-3 with each successive row counting up by 1. The counter works by ``passing'' values
along the glues of each tile. 

The south and east sides of each trit tile act as an input. The tile itself represents the result of
the operation on the input, in this case addition. For example, all tiles in category \verb|One| has
inputs that would add up to 1 in base three. Each trit category has to account for all possible
inputs that would add up to its corresponding value. The trit tiles are named according to its south
then east inputs.

The north and west sides of each trit tile work as outputs to give to the next tile. On the north
side, it passes its own value to the south side of the next trit tile. On the west side, it passes
along the carry value of the operation. E.g., if the inputs on the south and east sides were 2 and
1, the result of the trit is 0 and has to carry a 1 onto the next trit on the left.

The west glue of the \verb|RIGHT| tile always passes a value of 1 to the east side of a trit tile.
This results in adding a one to the value of the previous row of trit tiles since the trit
tiles of the previous row passes along its own values upwards. The bottom border, made up of
\verb|BOTTOM| tiles acts as a row with a value of 0 since each \verb|BOTTOM| tile passes along a 0
on its north side.

It's also important to see that all of the glue that passes along values are of strength-1 since the
system is temperature-2 and the trit tiles need two inputs.

\subsection*{Exercise 2} 
The goal of this temperature-2 system is to create a log-width binary counter i.e. produce a binary
counter that does not produce leading or trailing zeroes. In this system, there are 12 tiles evenly
split between two categories: \verb|Zero| and \verb|One| as shown in \cref{fig:exs-2}.

\begin{figure}[htbp]
    \centering
    \begin{subfigure}{\textwidth}
        \includesvg[width=.15\linewidth]{exercise_2-SEED}
        \includesvg[width=.15\linewidth]{exercise_2-00}
        \includesvg[width=.15\linewidth]{exercise_2-11}
        \includesvg[width=.15\linewidth]{exercise_2-1E}
        \includesvg[width=.15\linewidth]{exercise_2-11_G}
        \includesvg[width=.15\linewidth]{exercise_2-1E_G}
        \caption{\texttt{Zero} tiles.}
        \label{fig:exs-2_zero}
    \end{subfigure}
    \begin{subfigure}{\textwidth}
        \includesvg[width=.15\linewidth]{exercise_2-01}
        \includesvg[width=.15\linewidth]{exercise_2-10}
        \includesvg[width=.15\linewidth]{exercise_2-0E}
        \includesvg[width=.15\linewidth]{exercise_2-10_G}
        \includesvg[width=.15\linewidth]{exercise_2-0E_G}
        \includesvg[width=.15\linewidth]{exercise_2-E1_G}
        \caption{\texttt{One} tiles.}
        \label{fig:exs-2_one}
    \end{subfigure}
    \caption{Each category of tiles for exercise 2.}
    \label{fig:exs-2}
\end{figure}

This binary counter works similarly to the previous counter by thinking of its south and east sides
as inputs, its north and west sides as its outputs, the tile itself as the result of the operation
on the inputs, and the tiles passing values between each other through the glues. The seed aseembly
also contains a single \verb|SEED| tile.

In addition to passing values along the glues, it also passes an extra bit of information, $S$,
to indicate that it is the most significant bit of the current row, i.e. the leftmost column. This
is denoted by the presence of an \verb|S| in the glue or tile name. The tile knowing that it is the
most significant bit is important and will be discussed later.

Since this counter does not have a border on the right to give an input of value 1 on the east side
for each row, it is replaced by having a new input of value $\varepsilon$, denoted with an \verb|E|.
The least significant bit, i.e. the rightmost column, of each row always has an input of value
$\varepsilon$ on its east side. This means that the glues on the north and south sides of the least
significant bit has to be of strength-2, otherwise the system will never grow.

\subsection*{Exercise 3}
Make an assembly which grows to 100 rows tall and then stops growing. Each row must represent a
positive number in base-2. Each row can either be finite (like \#2) or have infinitely many leading
0's (like the example base-2 counter). In either case, the final assembly must contain exactly 100
rows of tiles. Your tile set for this problem should produce the same structure each time, no matter
what order individual tiles are added. Your tile set must work at temperature 2 for this problem.
Hint: not every row needs to represent a different value! This may seem tricky at first! 

\end{document}
